\begin{corollary}[无分辨率无关的常精度全息重构：噪声预算的指数收缩门槛]\label{cor:spg-dyadic-holographic-reconstruction-noise-budget}
沿用定理 \ref{thm:spg-dyadic-top-dimensional-holographic-inversion-exponential-ill-conditioning} 的设定。
设观测边界链为
\[
b=\partial_n x+\eta,
\]
其中 \(x\in C_n(\mathcal Q_m;\RR)\) 为真值、\(\eta\in C_{n-1}(\mathcal Q_m;\RR)\) 为 \(\ell^2\) 噪声。
对任意线性重构 \(x^\sharp:=Lb\)（\(L\partial_n=\mathrm{Id}\)），存在噪声实现使
\[
\norm{x^\sharp-x}_2\ \ge\ \frac{1}{\sqrt{2n}}\,2^{m/2}\,\norm{\eta}_2.
\]
特别地，若要求 \(\norm{x^\sharp-x}_2\le \delta\) 对所有满足 \(\norm{\eta}_2\le \varepsilon\) 的噪声均成立，则必须有
\[
\varepsilon\ \le\ \sqrt{2n}\,\delta\,2^{-m/2}.
\]
\leanverified{le\_div\_of\_pos\_mul\_le}
\leanverified{pos\_mul\_le\_iff\_le\_div}
\leanverified{paper\_noiseBudget\_threshold\_iff}
\leanverified{mul\_le\_of\_le\_div\_pos}
\leanverified{noiseBudget\_recovery\_symmetry}
\leanverified{noiseBudget\_contrapositive}
\leanverified{pos\_mul\_lt\_iff\_lt\_div}
\leanverified{paper\_noiseBudget\_algebra\_package}
\leanverified{paper\_spg\_dyadic\_holographic\_reconstruction\_noise\_budget}
因此当 \(m\to\infty\) 时，可容许的边界噪声幅度必须指数衰减；不存在对全部分辨率统一有效的常数精度全息重构。
\end{corollary}
\begin{proof}
由 \(x^\sharp-x=L(\partial_n x+\eta)-x=L\eta\) 得 \(\norm{x^\sharp-x}_2\le \norm{L}_{2\to 2}\norm{\eta}_2\)，并且存在某个单位向量方向的 \(\eta\) 使得 \(\norm{L\eta}_2=\norm{L}_{2\to 2}\norm{\eta}_2\)。
结合定理 \ref{thm:spg-dyadic-top-dimensional-holographic-inversion-exponential-ill-conditioning} 中 \(\norm{L}_{2\to 2}\ge \frac{1}{\sqrt{2n}}2^{m/2}\)，得到第一式下界。
若对所有 \(\norm{\eta}_2\le \varepsilon\) 都要求 \(\norm{x^\sharp-x}_2\le \delta\)，则必须有 \(\norm{L}_{2\to 2}\varepsilon\le \delta\)，代入上述下界即得所示必要条件。证毕。
\end{proof}

% --- Distillation writeback ---
\begin{proposition}[线性全息解码的精确填充半径门与双 lift 混淆]\label{distill:gromov-filling_budget_closure_families-spg-rad-linear-confusion-gate}
沿用推论 \ref{cor:spg-dyadic-holographic-reconstruction-noise-budget} 的 dyadic 链复形记号，并令
\[
L:C_{n-1}(\mathcal Q_m;\RR)\longrightarrow C_n(\mathcal Q_m;\RR),
\qquad L\partial_n=\mathrm{Id}.
\]
对 \(\delta>0\) 定义该线性左逆的填充半径门
\[
\operatorname{Rad}_L(\delta):=\sup\Bigl\{\varepsilon\ge0:
\|\eta\|_2\le\varepsilon\Longrightarrow \|L\eta\|_2\le\delta\Bigr\}.
\]
则
\[
\boxed{\ \operatorname{Rad}_L(\delta)=\frac{\delta}{\|L\|_{2\to2}}\ }.
\]
因此，给定 \(L\) 后，要求
\[
\|L(\partial_nx+\eta)-x\|_2\le\delta
\qquad\text{对所有 }x\text{ 与所有 }\|\eta\|_2\le\varepsilon\text{ 成立}
\]
当且仅当 \(\varepsilon\le \operatorname{Rad}_L(\delta)\)。此外定义两 lift 混淆门
\[
\operatorname{Amb}_{m,n}(\delta):=\frac12\inf\Bigl\{\|\partial_nh\|_2:\ h\in C_n(\mathcal Q_m;\RR),\ \|h\|_2>2\delta\Bigr\}.
\]
若 \(\varepsilon>\operatorname{Amb}_{m,n}(\delta)\)，则存在两个 bulk 链 \(x_0,x_1\) 与同一个观测边界 \(b\)，使得
\[
\|x_0-x_1\|_2>2\delta,\qquad
\|b-\partial_nx_0\|_2<\varepsilon,\qquad
\|b-\partial_nx_1\|_2<\varepsilon.
\]
于是任何解码器，不论线性与否，都不可能在全部 \(\varepsilon\)-噪声边界数据上保证误差 \(\le\delta\)。特别地，对任意线性左逆 \(L\)，推论 \ref{cor:spg-dyadic-holographic-reconstruction-noise-budget} 给出
\[
\operatorname{Rad}_L(\delta)\le \sqrt{2n}\,\delta\,2^{-m/2}.
\]
\end{proposition}
\begin{proof}
由于链空间有限维，算子范数 \(\|L\|_{2\to2}\) 在单位球面上取到。记 \(M:=\|L\|_{2\to2}\)。若 \(\varepsilon\le\delta/M\)，则对任意 \(\|\eta\|_2\le\varepsilon\)，有
\[
\|L\eta\|_2\le M\|\eta\|_2\le M\varepsilon\le\delta,
\]
故 \(\operatorname{Rad}_L(\delta)\ge\delta/M\)。反之，若 \(\varepsilon>\delta/M\)，取单位向量 \(u\) 满足 \(\|Lu\|_2=M\)，并取 \(t\) 使 \(\delta/M<t<\varepsilon\)。令 \(\eta=tu\)，则 \(\|\eta\|_2=t<\varepsilon\)，但
\[
\|L\eta\|_2=tM>\delta.
\]
因此任何大于 \(\delta/M\) 的半径都不满足定义中的一致误差要求，得到所示等式。

又因为 \(L\partial_n=\mathrm{Id}\)，对任意真值 \(x\) 与噪声 \(\eta\) 有
\[
L(\partial_nx+\eta)-x=L\eta.
\]
故 \(L\) 在 \(\varepsilon\)-噪声球上的 \(\delta\)-稳定性等价于 \(\|\eta\|_2\le\varepsilon\Rightarrow \|L\eta\|_2\le\delta\)，即等价于 \(\varepsilon\le\operatorname{Rad}_L(\delta)\)。

最后设 \(\varepsilon>\operatorname{Amb}_{m,n}(\delta)\)。由下确界定义，可取 \(h\in C_n(\mathcal Q_m;\RR)\) 满足
\[
\|h\|_2>2\delta,\qquad \frac12\|\partial_nh\|_2<\varepsilon.
\]
令
\[
x_0:=0,\qquad x_1:=h,\qquad b:=\frac12\partial_nh.
\]
则
\[
\|b-\partial_nx_0\|_2=\frac12\|\partial_nh\|_2<\varepsilon,\qquad
\|b-\partial_nx_1\|_2=\left\|\frac12\partial_nh-\partial_nh\right\|_2<\varepsilon,
\]
而 \(\|x_0-x_1\|_2=\|h\|_2>2\delta\)。若某解码器 \(D\) 对所有 \(\varepsilon\)-噪声边界均误差不超过 \(\delta\)，则同一个输入 \(b\) 必同时满足
\[
\|D(b)-x_0\|_2\le\delta,\qquad \|D(b)-x_1\|_2\le\delta,
\]
从而由三角不等式得 \(\|x_0-x_1\|_2\le2\delta\)，矛盾。最后的不等式直接由推论 \ref{cor:spg-dyadic-holographic-reconstruction-noise-budget} 中的 \(\|L\|_{2\to2}\ge(2n)^{-1/2}2^{m/2}\) 与 \(\operatorname{Rad}_L(\delta)=\delta/\|L\|_{2\to2}\) 相乘得到。证毕。
\end{proof}

% --- End distillation writeback ---

\endinput
